\section*{Bab \ref{ch:g11:genes-and-disease} --- Keragaman Genetik dan Penyakit}

\begin{solution}{exo:g11:genes-and-disease:1}
\omterm{def:g11:genes-and-disease:genetic}{Resesif}: penyakitnya muncul hanya dengan dua salinan \omterm{def:g10:universal-dna:gene}{alelnya}; \omterm{def:g11:genes-and-disease:genetic}{dominan}:
satu salinan sudah cukup. Seorang \omterm{def:g11:genes-and-disease:genetic}{pembawa} memuat satu \omterm{def:g10:universal-dna:gene}{alel} mutan yang
\omterm{def:g11:genes-and-disease:genetic}{resesif} dan tetap sehat.
\end{solution}

\begin{solution}{exo:g11:genes-and-disease:2}
Terkena $1/4$, \omterm{def:g11:genes-and-disease:genetic}{pembawa} $1/2$, bukan keduanya $1/4$.
\end{solution}

\begin{solution}{exo:g11:genes-and-disease:3}
I-1 $Aa$, I-2 $Aa$ (orang tua seorang anak yang terkena), II-2 $aa$,
III-2 $aa$.
\end{solution}

\begin{solution}{exo:g11:genes-and-disease:4}
Lelaki punya satu \omterm{def:g11:cell-cycle-mitosis:chromosome}{kromosom} X: \omterm{def:g10:universal-dna:gene}{alel} \omterm{def:g11:genes-and-disease:genetic}{resesif} padanya terungkapkan, karena
tidak ada salinan kedua yang menutupinya. Perempuan memerlukan \omterm{def:g10:universal-dna:gene}{alelnya}
pada kedua \omterm{def:g11:cell-cycle-mitosis:chromosome}{kromosom} X-nya.
\end{solution}

\begin{solution}{exo:g11:genes-and-disease:5}
Penyakit yang risikonya bergantung pada banyak \omterm{def:g10:universal-dna:gene}{alel} yang pengaruhnya
kecil dan pada lingkungannya: diabetes tipe 2, penyakit jantung (juga
asma, sebagian besar kanker).
\end{solution}

\begin{solution}{exo:g11:genes-and-disease:6}
$aa \times AA$: seluruh anaknya $Aa$, yaitu \omterm{def:g11:genes-and-disease:genetic}{pembawa} yang sehat. Anak
\omterm{def:g11:genes-and-disease:genetic}{pembawa} dengan seorang \omterm{def:g11:genes-and-disease:genetic}{pembawa}: $1/4$ terkena, $1/2$ \omterm{def:g11:genes-and-disease:genetic}{pembawa}, $1/4$
bukan \omterm{def:g11:genes-and-disease:genetic}{pembawa}.
\end{solution}

\begin{solution}{exo:g11:genes-and-disease:7}
II-1 adalah anak sehat dari dua \omterm{def:g11:genes-and-disease:genetic}{pembawa}: \omterm{def:g11:genes-and-disease:genetic}{pembawa} dengan peluang $2/3$.
Risikonya dengan seorang lelaki dari populasinya: $2/3 \times 1/25 \times 1/4
= 1/150$.
\end{solution}

\begin{solution}{exo:g11:genes-and-disease:8}
$1/2$ bagi tiap anaknya. Tidak satu pun dari empat: $(1/2)^4 = 1/16$.
\end{solution}

\begin{solution}{exo:g11:genes-and-disease:9}
$X^hY \times X^HX^H$: seluruh anak perempuannya $X^HX^h$, yaitu \omterm{def:g11:genes-and-disease:genetic}{pembawa}
yang sehat; seluruh anak lelakinya $X^HY$, sehat. Anak perempuan
\omterm{def:g11:genes-and-disease:genetic}{pembawa} dengan lelaki yang sehat: separuh anak lelakinya hemofilia,
separuh anak perempuannya \omterm{def:g11:genes-and-disease:genetic}{pembawa}.
\end{solution}

\begin{solution}{exo:g11:genes-and-disease:10}
Satu \omterm{def:g10:universal-dna:gene}{alelnya} memberi 400 dan 200; \omterm{def:g10:universal-dna:gene}{alel} yang lain, dengan delesi
tambahannya, memberi 300 dan 200. Pita pada 400, 300 dan 200 ---
tanpa 600.
\end{solution}

\begin{solution}{exo:g11:genes-and-disease:11}
Frekuensi sebuah \omterm{def:g10:universal-dna:gene}{alel} dalam sebuah populasi bergantung pada sejarahnya,
bukan hanya pada laju \omterm{def:g11:mutations:mutation}{mutasinya}: \omterm{def:g10:universal-dna:gene}{alel} Eropanya muncul lama berselang
lalu menyebar, barangkali karena \omterm{def:g11:genes-and-disease:genetic}{pembawanya} punya keuntungan tertentu
(ketahanan terhadap penyakit diare diduga demikian); di Asia timur hal
itu tidak terjadi.
\end{solution}

\begin{solution}{exo:g11:genes-and-disease:12}
Tiap anaknya adalah undian yang saling bebas: $1/4$ bagi yang kedua.
Sedikitnya satu dari dua berikutnya: $1 - (3/4)^2 = 7/16$. \omterm{def:g10:universal-dna:gene}{Alelnya}
dibagikan lagi dari awal pada tiap pembuahan; anak yang terkena
sebelumnya tidak menurunkan risiko anak berikutnya.
\end{solution}

\begin{solution}{exo:g11:genes-and-disease:13}
Sebelum ujinya, ia \omterm{def:g11:genes-and-disease:genetic}{pembawa} dengan peluang $1/25$; ujinya melewatkan
10\% \omterm{def:g11:genes-and-disease:genetic}{pembawanya}, sehingga di antara 1000 perempuan ada 40 \omterm{def:g11:genes-and-disease:genetic}{pembawa} yang
4 di antaranya diuji negatif, sedangkan 960 bukan \omterm{def:g11:genes-and-disease:genetic}{pembawa} diuji
negatif: peluang menjadi \omterm{def:g11:genes-and-disease:genetic}{pembawa} $4/964 \approx 1/240$. Risikonya
dengan \omterm{def:g11:genes-and-disease:genetic}{pembawa} yang diketahui: $1/240 \times 1/4
\approx 1/960$.
\end{solution}

\begin{solution}{exo:g11:genes-and-disease:14}
Dari banyak ke sedikit \omterm{def:g10:universal-dna:gene}{alel}, dalam keadaan kurang gerak: 34\% menjadi
9\%; dari kurang gerak ke aktif, dengan banyak \omterm{def:g10:universal-dna:gene}{alel}: 34\% menjadi 12\%.
Perubahan yang dapat dilakukan orang memberi penurunan yang nyaris
sama dengan perubahan yang tidak dapat dilakukannya. Skor risiko
berguna untuk mengenali siapa yang paling beroleh untung dari mengubah
cara hidupnya, bukan sebagai sebuah vonis.
\end{solution}

\begin{solution}{exo:g11:genes-and-disease:15}
Ujinya menawarkan kepastian --- kebebasan dari bayangan lima puluh
persen jika negatif, kemampuan merencanakan hidup dan keluarga jika
positif --- dan menuntut, jika positif, pengetahuan tentang penyakit
tak tersembuhkan puluhan tahun lebih awal, dengan pengaruh pada suasana
hati, asuransi dan kerabat yang berbagi \omterm{def:g10:universal-dna:gene}{gennya}. Tidak ada pengobatan
yang mengikuti dari hasilnya, sehingga satu-satunya alasan untuk tahu
adalah alasan orangnya sendiri; itulah sebabnya pilihannya, dan
saatnya, ada padanya.
\end{solution}

\begin{solution}{pb:g11:genes-and-disease:1}
\textbf{1.} \omterm{def:g11:genes-and-disease:genetic}{Resesif}: Zoé terkena dan kedua orang tuanya sehat,
sehingga masing-masing membawa satu \omterm{def:g10:universal-dna:gene}{alel} mutan yang tersembunyi.

\textbf{2.} Léa $Aa$, Tom $Aa$, Zoé $aa$.

\textbf{3.} $AA$ dengan peluang $1/3$, $Aa$ dengan $2/3$ (ia sehat,
sehingga $aa$ tersingkir).

\textbf{4.} Tom adalah \omterm{def:g11:genes-and-disease:genetic}{pembawa}, sehingga sedikitnya salah satu orang
tuanya juga; karena tidak ada orang yang terkena pada keturunan itu,
Ana menerima \omterm{def:g10:universal-dna:gene}{alelnya} dari orang tua tersebut dengan peluang $1/2$:
$1/2$ (dengan mengambil orang tua yang lain sebagai bukan \omterm{def:g11:genes-and-disease:genetic}{pembawa}).

\textbf{5.} $1/25$, yaitu angka populasinya.

\textbf{6.} $1/2 \times 1/25 \times 1/4 = 1/200$.

\textbf{7.} Terkena $1/4$; \omterm{def:g11:genes-and-disease:genetic}{pembawa} $1/2$.

\textbf{8.} $2/3 \times 1/25 \times 1/4 = 1/150$.

\textbf{9.} Zoé adalah $aa$; pasangannya $Aa$ dengan peluang $1/25$.
Terkena: $1/25 \times 1/2 = 1/50$. \omterm{def:g11:genes-and-disease:genetic}{Pembawa}: tiap anaknya menerima $a$
dari Zoé, sehingga ia \omterm{def:g11:genes-and-disease:genetic}{pembawa} dengan peluang $24/25 \times 1 + 1/25 \times
1/2 \approx 0.98$.

\textbf{10.} Lou adalah \omterm{def:g11:genes-and-disease:genetic}{pembawa} dengan peluang $2/3$; sepupunya
menerima $a$ dari Ana (\omterm{def:g11:genes-and-disease:genetic}{pembawa} dengan peluang $1/2$) dengan peluang
$1/2$, yaitu menjadi \omterm{def:g11:genes-and-disease:genetic}{pembawa} dengan peluang $1/4$ (dengan mengabaikan
Karim): $2/3 \times 1/4
\times 1/4 = 1/24$, empat kali lebih tinggi
daripada bagi Max, karena sepupunya berasal dari keluarga yang sama dan
jauh lebih mungkin daripada orang asing untuk membawa \omterm{def:g10:universal-dna:gene}{alel} yang sama.

\textbf{11.} Zoé: hanya 600 dan 300. Tom: 900, 600 dan 300. Bukan
\omterm{def:g11:genes-and-disease:genetic}{pembawa}: hanya 900.

\textbf{12.} $AA$: ia bukan \omterm{def:g11:genes-and-disease:genetic}{pembawa}; risiko bagi anaknya turun dari
$1/150$ menjadi nol (kecuali ada \omterm{def:g11:mutations:mutation}{mutasi} baru).

\textbf{13.} Ana adalah $Aa$: risikonya $1 \times 1/25 \times 1/4 = 1/100$.

\textbf{14.} Di antara 1000 orang seperti Karim, 40 adalah \omterm{def:g11:genes-and-disease:genetic}{pembawa} yang
4 di antaranya diuji negatif, dan 960 bukan \omterm{def:g11:genes-and-disease:genetic}{pembawa} diuji negatif: $4/964 \approx
1/240$. Risikonya: $1 \times 1/240 \times 1/4 \approx 1/960$.

\textbf{15.} $Aa$: seorang \omterm{def:g11:genes-and-disease:genetic}{pembawa} yang sehat, seperti ibunya.

\textbf{16.} $(1/25)^2 = 1/625$ pasangan suami istri adalah dua
\omterm{def:g11:genes-and-disease:genetic}{pembawa}; seperempat anaknya terkena: $1/2500$.

\textbf{17.} Persis frekuensi yang teramati: frekuensi \omterm{def:g11:genes-and-disease:genetic}{pembawanya} dan
aturan \omterm{def:g11:genes-and-disease:genetic}{resesifnya} menjelaskannya.

\textbf{18.} Dengan frekuensi \omterm{def:g11:genes-and-disease:genetic}{pembawa} $1/25$ dan frekuensi orang yang
terkena $1/2500$, \omterm{def:g10:universal-dna:gene}{alel} mutan yang dibawa \omterm{def:g11:genes-and-disease:genetic}{pembawa} yang sehat melampaui
\omterm{def:g10:universal-dna:gene}{alel} pada orang yang terkena dengan perbandingan $2 \times 1/25$
terhadap $2 \times
1/2500$, yaitu seratus berbanding satu. Kematian
orang yang terkena hanya menyingkirkan 1\% salinannya tiap keturunan;
sisanya diwariskan diam-diam oleh \omterm{def:g11:genes-and-disease:genetic}{pembawanya}.

\textbf{19.} Frekuensi \omterm{def:g10:universal-dna:gene}{alelnya} akan naik, tetapi sangat lambat: 1\%
salinan yang dahulu hilang tiap keturunan kini tersimpan, sehingga
perubahannya berorde 1\% frekuensinya tiap keturunan.

\textbf{20.} Dua \omterm{def:g11:genes-and-disease:genetic}{pembawa} berisiko mendapat satu anak yang terkena dari
empat; saudara sehat seorang anak yang terkena adalah \omterm{def:g11:genes-and-disease:genetic}{pembawa} dua kali
dari tiga; gelnya mengubah peluang menjadi \omterm{def:g11:enzymes-and-phenotype:phenotype}{genotipe} --- membersihkan
Max, membenarkan Ana, menenangkan tetapi tidak membersihkan Karim ---
dan mengenali janinnya sebagai \omterm{def:g11:genes-and-disease:genetic}{pembawa} yang sehat.
\end{solution}
